\documentclass[11pt]{article}
\usepackage[a4paper,margin=1in]{geometry}
\usepackage{fontspec}
\usepackage[boldfont=false]{xeCJK}
\setCJKmainfont{FandolSong-Regular.otf}
\usepackage{amsmath}
\usepackage{amssymb}
\usepackage{array}
\usepackage{booktabs}
\usepackage{graphicx}
\usepackage{adjustbox}
\usepackage{enumitem}
\setlist[itemize]{topsep=2pt,partopsep=0pt,parsep=0pt,itemsep=2pt,leftmargin=1.4em}
\setlist[enumerate]{topsep=2pt,partopsep=0pt,parsep=0pt,itemsep=2pt,leftmargin=1.6em}
\usepackage{parskip}
\usepackage{fvextra}
\fvset{breaklines=true,breakanywhere=true,fontsize=\small}
\setlength{\parindent}{0pt}

\begin{document}

\begin{center}
  {\LARGE\bfseries Gas exchange}\\[0.35em]
  {\large A Level Biology}
\end{center}
\vspace{0.6em}

\section*{The gas exchange system}

Your body needs to take in oxygen and get rid of carbon dioxide. This swap happens in the \textbf{gas exchange} 气体交换 system. Air follows this path into the body:

\begin{itemize}
\item down the \textbf{trachea} 气管 (the windpipe),

\item into two \textbf{bronchi} 支气管 (one to each lung),

\item into many smaller \textbf{bronchioles} 细支气管,

\item and finally into tiny air sacs called \textbf{alveoli} 肺泡, deep in the \textbf{lungs} 肺.

\end{itemize}

Each alveolus is wrapped in a network of \textbf{capillaries} 毛细血管, so air and blood are brought very close together.

\section*{The tissues of the airways and what they do}

\begin{center}
\adjustbox{max width=\linewidth}{%
\begin{tabular}{lll}
\toprule
\textbf{Tissue} & \textbf{Where it is} & \textbf{Function} \\
\midrule
\textbf{cartilage} 软骨 & C-shaped rings in the trachea and bronchi & holds the airway open so it cannot collapse when you breathe in \\
\textbf{ciliated epithelium} 纤毛上皮 & lining the trachea and bronchi & tiny hairs called \textbf{cilia} 纤毛 beat to sweep \textbf{mucus} 黏液 up towards the throat, away from the lungs \\
\textbf{goblet cells} 杯状细胞 and \textbf{mucous glands} 黏液腺 & in the lining of the airways & make the mucus, which traps dust and \textbf{microbes} 微生物 that you breathe in \\
\textbf{smooth muscle} 平滑肌 & in the walls of bronchi and bronchioles & contracts to make the airway narrower \\
\textbf{elastic fibres} 弹性纤维 & in the airway and alveolus walls & stretch when you breathe in, then spring back to help push air out \\
\textbf{squamous epithelium} 扁平上皮 & the very thin, flat lining of the alveoli & gives a very short distance for gases to cross \\
\bottomrule
\end{tabular}}
\end{center}

The cilia, goblet cells and mucous glands work together to keep the lungs clean and healthy: the mucus traps dirt and microbes, and the cilia carry it away to be swallowed.

\section*{Gas exchange in the alveoli}

The alveoli are excellent surfaces for exchanging gases, because they have:

\begin{itemize}
\item a very large total surface area (millions of tiny sacs),

\item very thin walls — the squamous epithelium of the alveolus and the capillary wall are each only one cell thick, so the distance to cross is tiny,

\item a rich blood supply from the capillary network,

\item a moist lining, so gases dissolve before crossing.

\end{itemize}

Gases move by \textbf{diffusion} 扩散 down their \textbf{concentration gradients} 浓度梯度:

\begin{itemize}
\item \textbf{oxygen} 氧气 is at a high concentration in the alveolar air and a low concentration in the blood, so it diffuses from the air into the blood.

\item \textbf{carbon dioxide} 二氧化碳 is at a high concentration in the blood and a low concentration in the alveolar air, so it diffuses from the blood into the air to be breathed out.

\end{itemize}

Breathing keeps fresh air in the alveoli, and the flowing blood keeps carrying gases away. Both of these keep the concentration gradients steep, so gas exchange stays fast.


\end{document}
